% Section 2.2.2 Analysis Results
% Owner: Slice 4 (issue #5).
\subsubsection{Analysis Results}
\label{sec:analysis-results}

We derive the following estimates from the quality requirements in the project
brief. These are workload and capacity calculations, not measured performance results.

\paragraph{Local availability.}
For an observation period \(T\), 99\% availability allows at most \(0.01T\) downtime.
Over an illustrative 30-day period, this is
\(30 \times 24 \times 60 \times 0.01 = 432\) minutes, or 7 hours 12 minutes.
Table~4.1 uses a 90-day period and lists 21 hours 36 minutes for 99\% availability
(Bass, Clements, and Kazman 53). Our observation period and downtime-accounting rules
remain to be agreed. Independently, the brief permits no individual local-inventory
outage longer than 30 minutes; a total budget does not override that cap.

\paragraph{Store buffer capacity.}
At 100 events/s, an illustrative 30-minute WAN interruption generates
\(100 \times 30 \times 60 = 180{,}000\) additional events per store.
The \Outbox needs capacity for that backlog, plus existing queued events, operating
headroom and storage overhead. A byte allocation requires the persisted event size.
The 30-minute WAN duration is a sizing assumption, not a maximum network outage
established by the brief's separate local-access limit.

\paragraph{Central throughput.}
Three stores produce up to \(3 \times 100 = 300\) new store events/s; ten produce
\(10 \times 100 = 1{,}000\). The \CFE and central processing path must also account
for \DC events, retries, subscriber fan-out and replay traffic. To drain a backlog,
effective processing capacity must exceed the ongoing arrival rate. We therefore
need workload validation with recovery traffic, not only a steady 1,000 events/s
input test; these calculations alone establish no achieved throughput.
